\documentclass[a4paper,11pt]{article}
\usepackage[utf8]{inputenc}
\usepackage[T1]{fontenc}
\usepackage[french]{babel}
\usepackage[left=1cm,right=1cm,top=.5cm,bottom=0cm]{geometry}
\usepackage{enumitem}
\usepackage{hyperref}
\usepackage{xcolor}
\usepackage{titlesec}
\usepackage{fontawesome5}
\usepackage{lmodern}
\usepackage[normalem]{ulem}

% --- Couleurs Startup / Tech (Violet-Indigo moderne) ---
\definecolor{primary}{RGB}{79, 70, 229}
\definecolor{secondary}{RGB}{80, 80, 80}

% --- Configuration des liens ---
\hypersetup{
    colorlinks=true,
    linkcolor=primary,
    filecolor=primary,
    urlcolor=primary,
}

% --- Formatage des titres ---
\titleformat{\section}{\large\bfseries\color{primary}\uppercase}{}{0em}{}[\titlerule]
\titlespacing{\section}{0pt}{8pt}{5pt}

% --- Commandes personnalisées ---
\newcommand{\resumeItem}[1]{\item\small{#1}}
\newcommand{\resumeSubheading}[4]{
  \vspace{-1pt}\item
    \begin{tabular*}{0.97\textwidth}[t]{l@{\extracolsep{\fill}}r}
      \textbf{#1} & \textit{\small #2} \\
      \textit{\small #3} & \textit{\small #4} \\
    \end{tabular*}\vspace{-5pt}
}

\begin{document}

% --- EN-TÊTE ---
\begin{center}
    {\Large \textbf{Loïc Aron MBASSI EWOLO}} \\ \vspace{4pt}
    \textbf{\large Alternance : Développeur IA \& Full-Stack – Agents Conversationnels} \\
    \textit{\small Élève Ingénieur en 2ème année, double diplôme ENSEA} \\ \vspace{4pt}
    \small
    \faMapMarker\ Cergy, France (Mobile Paris / Île-de-France) \quad
    \faPhone\ (+33) 7 59 52 33 98 \\ \vspace{2pt}
    \faEnvelope\ \href{mailto:loicaron.mbassiewolo@ensea.fr}{loicaron.mbassiewolo@ensea.fr} \quad
    \faLinkedin\ \href{https://www.linkedin.com/in/mbassi-loic-3942a9246/}{@mbassi-loic} \quad
    \faGithub\ \href{https://github.com/Nameless0l}{@Nameless0l} \quad
    \faGlobe\ \href{https://mbassiloic.tech}{mbassiloic.tech} \quad
    \faMedium\ \href{https://medium.com/@wwwmbassiloic}{medium}
\end{center}

\vspace{-8pt}

% --- PROFIL ---
\noindent\small{Élève-ingénieur en double diplôme à l'\textbf{ENSEA} avec une formation riche et variée en \textbf{Mathématiques, Électronique, Informatique et Programmation}. 
J'aime explorer de nouvelles approches et technologies ainsi que relever des défis open source. 
Ce dont témoignent la note de \textbf{19,25/20} en physique obtenue au Baccalauréat S, la 2\textsuperscript{e} place au hackathon IndabaX \textbf{Africa} et mes divers travaux en \textbf{IA conversationnelle}, \textbf{systèmes multi-agents} et développement \textbf{Full-Stack PHP/JavaScript}.}

\vspace{4pt}

% --- FORMATION ---
\section{Formation}
\begin{itemize}[leftmargin=*, label={.}, nosep]
    \item \textbf{Double diplôme : École Nationale Supérieure de l'Électronique et de ses Applications (ENSEA)} \hfill \textit{2025 -- Présent} \\
    \small{2ème année - Spécialité Intelligence Artificielle et Traitement du Signal, Ingénieur de conception}
    \vspace{2pt}
    \item \textbf{École Nationale Supérieure Polytechnique de Yaoundé — Génie Informatique} \hfill \textit{2021 -- Présent} \\
    \small{Niveaux 4 \& 5 : \textbf{Systèmes multi-agents}, NLP, conception logicielle, architectures distribuées} \\
    \small{Niveaux 1 à 3 : Logique, algorithmique, programmation C, POO Java | \textbf{Licence 1} Informatique en parallèle}
\end{itemize}

% --- COMPÉTENCES TECHNIQUES ---
\section{Compétences Techniques}
\begin{itemize}[leftmargin=*, nosep]
    \item \small{\textbf{IA Conversationnelle \& Agents :} Google ADK, LangChain, RAG, Prompt Engineering, APIs OpenAI/Gemini, NLP, BERT}
    \item \small{\textbf{Backend :} \textbf{PHP/Laravel}, Python, Java/Spring Boot | APIs RESTful, Microservices, RabbitMQ}
    \item \small{\textbf{Frontend :} \textbf{JavaScript}, React.js, Next.js, TypeScript, HTML/CSS}
    \item \small{\textbf{Outils \& DevOps :} Git/GitHub, Docker, MySQL, PostgreSQL, MongoDB, Redis | CI/CD, Tests (TDD)}
    \item \small{\textbf{IA/ML :} PyTorch, TensorFlow, Hugging Face Transformers, Computer Vision, Fine-tuning LLM}
\end{itemize}

% --- PROJETS IA CONVERSATIONNELLE & AGENTS ---
\section{Projets – IA Conversationnelle \& Systèmes Multi-Agents}
\begin{itemize}[leftmargin=*, label={}]

    \item \textbf{Système Multi-Agents avec Google ADK — Chatbot Agricole} \hfill \textit{2024 -- 2025, personnel}
    \vspace{-2pt}
    \begin{itemize}[leftmargin=0.4cm, nosep, label=\tiny$\bullet$]
        \item \small{Développement d'un système d'agents IA collaboratifs orchestrant \textbf{5 agents spécialisés} pour fournir dans un \textbf{chatbot} des conseils personnalisés aux agriculteurs camerounais}
        \item \small{Architecture de conversation avec gestion des scénarios, objections et propositions contextuelles}
        \item \small{Orchestration multi-agents avec Google ADK, \textbf{RAG} (Retrieval-Augmented Generation), intégration API Gemini}
        \item \small{\textbf{Stack :} Google ADK, Gemini 2.0 Flash, FastAPI, Docker, LangChain, Poetry — vectorisation de données}
    \end{itemize}
    \vspace{4pt}

    \item \textbf{Laravel API Generator — Génération de Code Assistée par LLM} \hfill \textit{2024 -- Présent, Open Source}
    \vspace{-2pt}
    \begin{itemize}[leftmargin=0.4cm, nosep, label=\tiny$\bullet$]
        \item \small{Conception d'un package assisté par IA pour la génération de projets (à partir de diagrammes UML) en suivant les bonnes pratiques}
        \item \small{Intégration de \textbf{LLM} pour l'assistance contextuelle au code et prompt engineering}
        \item \small{Parsing XML, métaprogrammation Python, tests unitaires automatisés, documentation auto — \href{https://github.com/Nameless0l}{\faGithub}}
    \end{itemize}
    \vspace{4pt}

    \item \textbf{Women Safe — Détection de Contenus Sexistes (NLP)} \hfill \textit{2023, MLPC Competition}
    \vspace{-2pt}
    \begin{itemize}[leftmargin=0.4cm, nosep, label=\tiny$\bullet$]
        \item \small{Développement d'un modèle de détection d'expressions sexistes avec \textbf{BERT} et fine-tuning}
        \item \small{Analyse de texte en français, gestion des biais, évaluation éthique — 2\textsuperscript{e} place compétition}
    \end{itemize}

\end{itemize}

% --- EXPÉRIENCE PROFESSIONNELLE ---
\section{Expériences Professionnelles — Développement Full-Stack}
\begin{itemize}[leftmargin=*, label={}]

    \resumeSubheading
      {Développeur Full-Stack (Fintech) — CSC}{Oct. 2024 -- Août 2025}
      {Douala, Cameroun}{}
      \resumeItem{Conception et développement d'une API fintech sécurisée et d'interfaces de gestion}
      \begin{itemize}[leftmargin=0.4cm, nosep, label=\tiny$\bullet$]
        \item \small{Architecture backend sécurisée (\textbf{Laravel}/Docker, Redis), APIs RESTful, gestion des accès concurrents}
        \item \small{Tests unitaires \& d'intégration (TDD), revue de code, CI/CD, méthodologie \textbf{Agile/Scrum}}
      \end{itemize}
    \vspace{2pt}

    \resumeSubheading
      {Développeur Full-Stack — SOSDOCTA, Belgique}{Fév. 2022 -- Nov. 2024}
      {Remote — Plateforme de télémédecine}{}
      \resumeItem{Application de consultation et gestion des patients avec télémédecine intégrée}
      \begin{itemize}[leftmargin=0.4cm, nosep, label=\tiny$\bullet$]
        \item \small{\textbf{Laravel}, MySQL, JavaScript — Dashboard client, APIs de paiement, notifications}
        \item \small{Maintenance 2+ ans, gestion de projet en équipe distribuée, conformité RGPD}
      \end{itemize}
\end{itemize}

% --- AUTRES PROJETS ---
\section{Autres Projets Significatifs}
\begin{itemize}[leftmargin=*, label={}]

    \item \textbf{Nettoie-Québec, Canada} \hfill \textit{2024, freelance}
    \vspace{-2pt}
    \begin{itemize}[leftmargin=0.4cm, nosep, label=\tiny$\bullet$]
        \item \small{Application web de réservation en ligne — Laravel, JavaScript, système de notifications email}
    \end{itemize}
    \vspace{2pt}

    % \item \textbf{YembaAudio — Reconnaissance Vocale Langue Locale} \hfill \textit{2025, personnel}
    % \vspace{-2pt}
    % \begin{itemize}[leftmargin=0.4cm, nosep, label=\tiny$\bullet$]
    %     \item \small{Système de reconnaissance vocale pour le Yemba (NLP), Next.js, TensorFlow, transfer learning}
    % \end{itemize}
    % \vspace{2pt}

    \item \textbf{GoFind — Architecture Microservices} \hfill \textit{2024, académique}
    \vspace{-2pt}
    \begin{itemize}[leftmargin=0.4cm, nosep, label=\tiny$\bullet$]
        \item \small{Plateforme microservices : Laravel, NestJS, MongoDB, communication inter-services avec \textbf{RabbitMQ}}
    \end{itemize}

\end{itemize}

% --- RECHERCHE ---
\section{Expérience en Recherche}
\begin{itemize}[leftmargin=*, label={}]
    \resumeSubheading
      {Stage de Recherche à distance — MILA}{Juin -- Août 2025}
      {Institut Québécois d'Intelligence Artificielle}{Québec, Canada}
      \resumeItem{Grokking (de OpenAI), Généralisation des réseaux de neurones après sur-apprentissage}
      \begin{itemize}[leftmargin=0.4cm, nosep, label=\tiny$\bullet$]
        \item \small{Reproduction d'expériences SOTA, implémentation de pipelines de tests sous \textbf{PyTorch}}
      \end{itemize}
\end{itemize}

% --- HACKATHONS & LEADERSHIP ---
\section{Hackathons, Leadership \& Certifications}
\begin{itemize}[leftmargin=*, nosep]
    \item \small{\textbf{Hackathons :} 2\textsuperscript{e} place IndabaX Africa (2025) | 1\textsuperscript{re} place CITS National \& Mountain Hub Regional (2024)}
    \item \small{\textbf{Certifications :} Supervised ML (DeepLearning.AI), Python (IBM), Agents IA (IBM, en cours)}
    \item \small{\textbf{Langues :} Français (Langue maternelle), Anglais (Communication scientifique)}
\end{itemize}

\end{document}
